\documentclass[12pt,a4paper]{article}

\usepackage[margin=1in]{geometry}
\usepackage[T1]{fontenc}
\usepackage{lmodern}
\usepackage{microtype}
\usepackage{xcolor}
\usepackage{graphicx}
\usepackage{booktabs}
\usepackage{enumitem}
\usepackage{listings}
\usepackage{caption}
\usepackage{float}
\usepackage[hidelinks]{hyperref}

\definecolor{codegray}{RGB}{246,246,246}
\definecolor{commentgreen}{RGB}{30,110,45}
\definecolor{keywordblue}{RGB}{0,55,135}

\lstdefinestyle{cstyle}{
    language=C,
    backgroundcolor=\color{codegray},
    basicstyle=\ttfamily\footnotesize,
    keywordstyle=\color{keywordblue}\bfseries,
    commentstyle=\color{commentgreen},
    numbers=left,
    numberstyle=\tiny\color{gray},
    numbersep=8pt,
    frame=single,
    breaklines=true,
    keepspaces=true,
    showstringspaces=false,
    captionpos=b
}

\lstdefinestyle{terminal}{
    backgroundcolor=\color{codegray},
    basicstyle=\ttfamily\footnotesize,
    frame=single,
    breaklines=true,
    keepspaces=true,
    columns=fullflexible,
    captionpos=b
}

\title{\textbf{Problem Session 4}\\
Memory Class and API in C Programming}
\author{%
    \begin{tabular}{ll}
        Kiatisak Markmeeshap & 67070501005\\
        Tithinan Sobking & 67070501018\\
        Worawut Sereethai & 67070501040\\
        Chanya Poolketkij & 67070501058
    \end{tabular}%
}
\date{}

\begin{document}

\maketitle
\tableofcontents
\newpage

\section{Environment and Compilation}

All experiments were performed in Ubuntu running under WSL. The source files,
executables, and output logs were kept in the same working directory.

\begin{lstlisting}[style=terminal,caption={Working directory and compilation examples}]
cd ~/cpe333/ps4
gcc -Wall -Wextra -std=c11 -O0 -g -o static_lab static_lab.c
gcc -Wall -Wextra -std=c11 -O0 -g -o static_lab_nopie static_lab.c -no-pie
./static_lab
./static_lab_nopie
\end{lstlisting}

The programs use standard C entry points with \texttt{int main(void)}. Pointer
values are cast to \texttt{void *} when printed with the \texttt{\%p} format
specifier.

\section{Static Storage Class}

\subsection{Original program}

The original program declares \texttt{y} as a static local variable. It is
initialized once and keeps its value between loop iterations.

\lstinputlisting[style=cstyle,caption={Static storage class program}]
{../static_lab.c}

\subsection{Results from three runs}

The program was compiled and executed three times without \texttt{-no-pie}.
The address of \texttt{y} stays constant within one execution, while its
address can change between executions because of position-independent loading
and address randomization.

\begin{lstlisting}[style=terminal,caption={Selected output from the original static program}]
The value of y is 6
The address of y is 0x5eb92dd39010

The value of y is 7
The address of y is 0x5eb92dd39010

The value of y is 8
The address of y is 0x5eb92dd39010
\end{lstlisting}

The complete three-run output is included in Appendix~\ref{app:outputs}.

\subsection{Program after removing \texttt{static}}

After removing the keyword, \texttt{y} becomes an automatic local variable.
It is initialized to 5 again on every loop iteration, so the printed value is
6 on every iteration.

\lstinputlisting[style=cstyle,caption={Static keyword removed}]
{../static_lab_no_static.c}

\begin{lstlisting}[style=terminal,caption={Output after removing static}]
The value of y is 6
The address of y is 0x7ffde7b84c40

The value of y is 6
The address of y is 0x7ffde7b84c40

The value of y is 6
The address of y is 0x7ffde7b84c40
\end{lstlisting}

The address remains the same during this execution because the same stack
location is reused by the loop. It is not guaranteed to be the same across
separate executions.

\subsection{Comparison with \texttt{-no-pie}}

With the static program compiled using \texttt{-no-pie}, the static variable
was located at the same address in all three runs in this environment:
\texttt{0x404018}. The values still followed the same sequence 6, 7, and 8.
For the non-static program, the stack address still changed between runs,
because disabling PIE does not disable stack address randomization.

\section{Extern Storage Class}

\subsection{Original program}

The global variable \texttt{x} is defined once and referenced from both
\texttt{main()} and \texttt{display()} using \texttt{extern}. Both functions
therefore access the same object.

\lstinputlisting[style=cstyle,caption={Extern storage class program}]
{../extern_demo.c}

\begin{lstlisting}[style=terminal,caption={Original extern output}]
Print value of x: 20
Address of x (in main function): 0x5bb79ed55010

Display value of x: 20
Address of x (in display function): 0x5bb79ed55010
\end{lstlisting}

The value and address are identical in both functions. With the default PIE
build, the global address may change between executions. With \texttt{-no-pie},
the global address was stable at \texttt{0x404018} in this environment.

\subsection{Removing \texttt{extern} in \texttt{main()}}

The modified line becomes \texttt{int x;}. This declares a new local variable
that hides the global variable. It is not initialized, so reading it produces
an indeterminate value and the compiler warns about using an uninitialized
variable. The \texttt{display()} function still accesses the global value 20.

\lstinputlisting[style=cstyle,caption={Extern removed from main}]
{../extern_demo_no_extern_main.c}

\begin{lstlisting}[style=terminal,caption={Output after removing extern in main}]
Print value of x: 32767
Address of x (in main function): 0x7fff7d990da4

Display value of x: 20
Address of x (in display function): 0x5ca88dea4010
\end{lstlisting}

The exact value printed by \texttt{main()} is not predictable and may differ
between runs. The important observation is that the two functions now refer to
different objects.

\section{Memory Allocation}

\subsection{Original program}

The program creates an integer pointer \texttt{a}, a floating-point pointer
\texttt{b}, and a fixed-size array \texttt{c}. It then allocates memory for
\texttt{a} using \texttt{malloc()} and expands it using \texttt{realloc()}.

\lstinputlisting[style=cstyle,caption={Memory allocation program}]
{../memory_demo.c}

\subsection{Meaning of the addresses}

\begin{itemize}[leftmargin=*]
    \item \texttt{\&a} and \texttt{\&b} are the addresses of the pointer variables themselves.
    \item \texttt{a} and \texttt{b} are the addresses stored inside the pointers.
    \item After \texttt{malloc()}, \texttt{a} points to a block on the heap.
    \item \texttt{c} and \texttt{\&c[0]} have the same base address because an array name decays to a pointer to its first element.
    \item \texttt{\&a[9]} and \texttt{\&c[9]} are offset by nine integer elements from the first element.
    \item \texttt{realloc()} may keep the original block or move it to a new location if a larger contiguous block is required.
\end{itemize}

\begin{lstlisting}[style=terminal,caption={Memory allocation output before enabling b}]
Address of Pointer
>>0x7ffd5e8f0be0
>>0x7ffd5e8f0be8

Effective Address
>>(nil)
>>(nil)

After malloc Pointer a
>>0x6467b2cec2b0
>>0x6467b2cec2b0
>>0x6467b2cec2d4

Array c
>>0x7ffd5e8f0bf0
>>0x7ffd5e8f0bf0
>>0x7ffd5e8f0c14

After realloc Pointer a
>>0x6467b2cec2b0
>>0x6467b2cec2b0
>>0x6467b2cec2d4
>>0x6467b2ced24c
\end{lstlisting}

The initial effective addresses are not valid allocated objects because the
pointers have not been initialized. Their printed values are therefore not
portable or predictable. The observed \texttt{(nil)} output is only one result
from this execution.

\subsection{After uncommenting the allocation of b}

The second version allocates ten floating-point elements for \texttt{b} before
calling \texttt{realloc()} on \texttt{a}.

\lstinputlisting[style=cstyle,caption={Memory allocation with b enabled}]
{../memory_demo_with_b.c}

\begin{lstlisting}[style=terminal,caption={Memory allocation output after enabling b}]
After malloc Pointer b
>>0x57232be552e0
>>0x57232be552e0
>>0x57232be55304

After realloc Pointer a
>>0x57232be55310
>>0x57232be55310
>>0x57232be55334
>>0x57232be562ac
\end{lstlisting}

In this run, allocating \texttt{b} occupied space immediately after the first
allocation for \texttt{a}. When \texttt{a} was expanded, the allocator moved it
to another address. This behavior is implementation-dependent, so another run
or another machine may produce different addresses.

\section{Discussion}

The experiments show that storage class affects lifetime, linkage, and the
location used by a variable. A static local variable retains its value after a
function block is left, while an automatic variable is initialized again on
each entry to the block. The \texttt{extern} declaration allows code in
different functions or source files to refer to one global object.

The address experiments also show the effect of process memory layout. PIE
allows the executable's data to be relocated between runs, while \texttt{-no-pie}
makes the executable's fixed data address stable in this environment. Heap
addresses remain controlled by the memory allocator and can change when
\texttt{malloc()} or \texttt{realloc()} changes the allocation layout.

\section{Conclusion}

Static storage preserves a variable's value for the lifetime of the program.
Extern storage provides access to a global object from another scope. Finally,
\texttt{malloc()} and \texttt{realloc()} demonstrate dynamic heap allocation,
while the address of an array declared inside a function belongs to the stack.

\appendix
\section{Complete Experiment Outputs}\label{app:outputs}

\lstinputlisting[style=terminal,caption={Complete static experiment output}]
{../outputs/static_runs.txt}

\lstinputlisting[style=terminal,caption={Complete extern experiment output}]
{../outputs/extern_runs.txt}

\lstinputlisting[style=terminal,caption={Complete memory experiment output}]
{../outputs/memory_runs.txt}

\end{document}
